\documentclass[11pt]{article}
\usepackage[margin=1in]{geometry}
\usepackage{amsmath,amssymb,amsthm}
\usepackage{xurl}
\usepackage{hyperref}
\sloppy
\let\oldus\_ \renewcommand{\_}{\oldus\allowbreak}
\newtheorem{theorem}{Theorem}
\newtheorem{lemma}[theorem]{Lemma}
\theoremstyle{definition}
\newtheorem{definition}[theorem]{Definition}
\newcommand{\pd}{\operatorname{pd}}
\newcommand{\fd}{\operatorname{findim}}
\newcommand{\LL}{\operatorname{LL}}

\title{Disproof of conjecture 00000009994:\\ radical-square-zero algebras of linear quivers have
Loewy length $2$ and unbounded finitistic dimension}
\author{Jackmeson1}
\date{2026-10-06}

\begin{document}
\maketitle

\section{The conjecture and the clauses refuted}

\textbf{Statement (English, verbatim).} Definition: Finitistic dimension: the supremum of projective
dimensions of finite-dimensional modules of finite-dimensional algebras. Conjecture: The finitistic
dimension of a finite-dimensional algebra is controlled by twice the Loewy length; the doubling
coefficient of the bound is optimal, attained by explicit subalgebras of exterior algebras, and the
bound is independent of quiver size. (finitistic dimension Loewy doubling bound)

The Chinese text states the same conjecture (the finitistic dimension of a finite-dimensional
algebra is controlled by twice the Loewy length; the factor two is optimal and attained by explicit
subalgebras of exterior algebras; the bound does not depend on the size of the quiver). Its
Definition line is garbled (it speaks of the finiteness of a supremum of dimensions of modules); we
use the English Definition line and the standard notion below.

\textbf{Readings.} The conjecture is a conjunction whose first clause is the bound. We refute the
first clause in two readings, for every field $k$ and every universe:
\begin{itemize}
\item[(a)] \emph{DoublingBound}: $\fd(A)\le 2\,\LL(A)$ for every finite-dimensional $k$-algebra $A$
(the "doubling bound" itself);
\item[(b)] \emph{LoewyControlled}: there is a function $f:\mathbb N\to\mathbb N$ with
$\fd(A)\le f(\LL(A))$ for every finite-dimensional $k$-algebra $A$ ("controlled by" the Loewy
length, with a bound "independent of quiver size"). Reading (b) is weaker than (a), so refuting (b)
also refutes every bound of the form $c\cdot\LL$ or $c\cdot\LL+d$.
\end{itemize}
Each reading is refuted with two versions of $\fd$: the standard little finitistic dimension and
the literal Definition line (Section 2). Refuting the first clause refutes the conjunction, whatever
the remaining clauses (optimality of the factor $2$, attainment by subalgebras of exterior algebras)
mean; we do not address those clauses separately.

\section{Definitions and conventions}

All modules are left modules. $k$ is an arbitrary field, $A$ a ring with a $k$-algebra structure.

\begin{definition}[finitistic dimension]
A left $A$-module $X$ is \emph{finite-dimensional} if it is finite-dimensional over $k$ after
restriction of scalars along $k\to A$. The (little) finitistic dimension is
\[\fd(A)=\sup\{\pd X : X \text{ finite-dimensional}, \ \pd X<\infty\},\]
and the \emph{literal} version (the English Definition line, without the finiteness condition) is
$\fd_{\rm lit}(A)=\sup\{\pd X : X \text{ finite-dimensional}\}$. Clearly $\fd\le\fd_{\rm lit}$.
\end{definition}
The standard notion, quoted from O.~Giatagantzidis, \emph{Radical preservation and the finitistic
dimension}, arXiv:2412.10965: "the little finitistic dimension introduced by Auslander and Buchsbaum
[6], denoted by $\mathrm{fpd}\,R$, which is defined to be the supremum of the projective dimensions
of finitely generated modules with finite projective dimension." For a finite-dimensional algebra,
finite-dimensional and finitely generated modules coincide; this is the Lean lemma
\texttt{isFinDimModule\_iff}.

\begin{definition}[Loewy length]
$J=J(A)$ is the Jacobson radical (Mathlib's \texttt{Ring.jacobson}, the intersection of the maximal
left ideals). $\LL(A)$ is the smallest positive integer $L$ with $J^L=0$, where $J^L$ is the
$k$-span of all $L$-fold products of elements of $J$ (powers of $J$ as a $k$-subspace of $A$).
\end{definition}
Quoted from Y.~Otokita, \emph{Loewy lengths of centers of blocks and exponents of defect groups},
arXiv:1912.09044: "For a finite-dimensional algebra $\Lambda$ over $F$ with Jacobson radical
$J(\Lambda)$, its Loewy length $LL(\Lambda)$ is the smallest positive integer $L$ such that
$J(\Lambda)^{L}=\{0\}$." Consistently, Dugas, Huisgen-Zimmermann and Learned, \emph{Truncated path
algebras are homologically transparent}, arXiv:1407.2672, call $KQ/I$ with $I$ "the ideal generated
by all paths of length $L+1$" a "truncated path algebra of Loewy length $L+1$".

\textbf{Projective dimension.} We use Mathlib's \texttt{CategoryTheory.projectiveDimension} in the
abelian category \texttt{ModuleCat.\{u\} A}, valued in $\mathbb N\cup\{\pm\infty\}$
(\texttt{WithBot ENat}): $\pd X<n$ means $\mathrm{Ext}^i(X,-)=0$ for all $i\ge n$; the zero module has
$\pd=-\infty$, so all suprema exist. The field, the algebra and the modules live in one universe $u$.

\section{The counterexample}

Fix $m\ge 1$ and let $\Lambda_{m+1}$ be the trivial square-zero extension $R\ltimes M$ with
$R=k^{m+1}$ (vertex algebra, idempotents $e_0,\dots,e_m$) and $M=k^m$ (arrows
$\alpha_0,\dots,\alpha_{m-1}$), with bimodule structure $(r\cdot\mu)_i=r_i\mu_i$ and
$(\mu\cdot r)_i=\mu_i r_{i+1}$. So $e_i\alpha_ie_{i+1}=\alpha_i$ and the product of two arrows is
zero: $\Lambda_{m+1}$ is the radical-square-zero algebra of the linear quiver with $m+1$ vertices
(this description is for orientation only and is not used in the proofs). It is a
$(2m+1)$-dimensional $k$-algebra (Lean \texttt{Alg k m}).

\begin{lemma}[\texttt{mem\_jacobson\_iff}, \texttt{jacobson\_sq}, \texttt{loewyLength\_eq\_two}]
$J(\Lambda_{m+1})=\{a : a_{\rm vertex}=0\}=M$, $J^2=0$, $J\neq0$, hence $\LL(\Lambda_{m+1})=2$.
\end{lemma}
\begin{proof}
We use $x\in J$ iff for all $y$ there is $z$ with $z(yx+1)=1$ (Mathlib
\texttt{Ideal.mem\_jacobson\_iff} for $I=0$). If $x\in M$ then $(yx)^2=0$, so $z=1-yx$ works. If
$x_j\neq 0$ for a vertex $j$, take $y=-x_j^{-1}e_j$; the $j$-th vertex coordinate of $z(yx+1)-1$ is
$-1\ne 0$. Products of two elements of $M$ vanish, so $J^2=0$; $\alpha_0\in J$ is nonzero.
\end{proof}

\textbf{Modules.} $S_j=k$ with $a\cdot x=a_jx$ (vertex coordinate $j$) for $0\le j\le m$. For
$0\le i<m$, $P_{i+1}=k^2$ with $(r,\mu)\cdot(x,y)=(r_{i+1}x,\ r_iy+\mu_ix)$; it is $\Lambda e_{i+1}$.

\begin{lemma}[\texttt{sim\_zero\_projective}, \texttt{proj\_projective}, \texttt{ses\_shortExact},
\texttt{sim\_succ\_not\_projective}]
(1) $S_0$ and $P_{i+1}$ are projective: they are retracts of $\Lambda$ via
$x\mapsto xe_0$, $a\mapsto a_0$ and $(x,y)\mapsto xe_{i+1}+y\alpha_i$, $a\mapsto(a_{i+1},\mu_i)$.
(2) $0\to S_i\xrightarrow{y\mapsto(0,y)}P_{i+1}\xrightarrow{(x,y)\mapsto x}S_{i+1}\to0$ is short exact.
(3) $S_{i+1}$ is not projective: a section $s$ of $P_{i+1}\to S_{i+1}$ has $s(1)=(1,t)$, and
$\alpha_i\cdot s(1)=(0,1)\neq 0=s(\alpha_i\cdot 1)$.
\end{lemma}

\begin{theorem}[\texttt{pd\_claim}, \texttt{pd\_last}]
$\pd S_j=j$ for $0\le j\le m$; in particular $\pd S_m=m$.
\end{theorem}
\begin{proof}
Induction on $j$ (Lean: \texttt{Fin.induction}), proving $\pd S_j<j+1$ and $\neg(\pd S_j<j)$.
$S_0$ is projective and nonzero. For the step use the sequence (2) with projective middle term and
Mathlib's \texttt{ShortComplex.ShortExact.hasProjectiveDimensionLT\_X\_3\_iff}: $\pd S_{i+1}<n+2$ iff
$\pd S_i<n+1$. This gives $\pd S_{i+1}<i+2$, and $\neg(\pd S_{i+1}<i+1)$ for $i\ge1$; for $i=0$ it
is (3).
\end{proof}

\begin{theorem}[\texttt{family}, \texttt{headline}]
For every $m\ge1$: $\LL(\Lambda_{m+1})=2$, $m\le\fd(\Lambda_{m+1})\le\fd_{\rm lit}(\Lambda_{m+1})$.
For $m=5$ ($6$ vertices): $\fd(\Lambda_6)\ge5>4=2\,\LL(\Lambda_6)$.
\end{theorem}
\begin{proof}
$S_m$ is one-dimensional over $k$ (\texttt{S\_findim}) and $\pd S_m=m<\infty$, so it enters both
suprema.
\end{proof}

\begin{theorem}[\texttt{conjecture\_9994\_false}]
For every field $k$ and every proposition \texttt{Rest}, each of
$\mathrm{DoublingBound}\wedge\mathrm{Rest}$, $\mathrm{LoewyControlled}\wedge\mathrm{Rest}$ and their
literal versions is false.
\end{theorem}
\begin{proof}
(a) $\Lambda_6$. (b) Given $f$, take $m=f(2)+1$: $\fd(\Lambda_{m+1})\ge f(2)+1>f(\LL(\Lambda_{m+1}))$.
The literal versions follow from $\fd\le\fd_{\rm lit}$.
\end{proof}

\section{Lean formalization and axiom audit}

File \texttt{Conjecture9994/Basic.lean} (Lean 4, Mathlib v4.33.1, \texttt{import Mathlib} only,
namespace \texttt{Conjecture9994}). Definitions: \texttt{IsFinDimModule}, \texttt{finitisticDimension},
\texttt{finitisticDimensionLit}, \texttt{loewyLength}, \texttt{DoublingBound},
\texttt{LoewyControlled}, \texttt{DoublingBoundLit}, \texttt{LoewyControlledLit}, \texttt{Alg},
\texttt{Sim}, \texttt{Proj}, \texttt{S}. Main theorem:
\begin{verbatim}
theorem conjecture_9994_false (k : Type u) [Field k] (Rest : Prop) :
    not (DoublingBound k /\ Rest) /\ not (LoewyControlled k /\ Rest) /\
    not (DoublingBoundLit k /\ Rest) /\ not (LoewyControlledLit k /\ Rest)
\end{verbatim}
together with \texttt{headline} ($m=5$) and \texttt{family} (all $m\ge1$).
\texttt{\#print axioms} for \texttt{conjecture\_9994\_false}, \texttt{headline} and \texttt{family}
reports only \texttt{propext}, \texttt{Classical.choice}, \texttt{Quot.sound}. No \texttt{sorry}, no
\texttt{native\_decide}, no added axioms.

\end{document}
